\subsection{فعالیت بالا رفتن از جعبه}

\subsubsection{هدف از افزودن این فعالیت}

فعالیت بالا رفتن از جعبه یا پله، از نظر عملکردی به توان اندام‌های تحتانی، کنترل تعادل، هماهنگی تنه و انتقال وزن وابسته است. این فعالیت نسبت به ایستادن پیچیده‌تر و نسبت به پریدن کنترل‌شده‌تر است و می‌تواند اطلاعات ارزشمندی درباره توانایی انجام حرکات عملکردی روزمره فراهم کند. در افراد دارای ضعف عضلانی، اجرای این فعالیت ممکن است با کندی حرکت، استفاده از حرکات جبرانی، تغییر الگوی قرارگیری پا یا ناپایداری هنگام بالا رفتن همراه شود.

در نسخه فعلی، خروجی‌های کامل عددی و تصویری این فعالیت هنوز وارد فصل نشده‌اند. این بخش به عنوان قالب آماده برای زمانی نوشته شده است که داده‌های این فعالیت ثبت و خروجی‌های مدل تولید شوند.

% TODO: پس از تولید خروجی‌ها، شکل‌های این فعالیت را در مسیر Report/Images/Generated/TaskResults/Box_Step یا نام نهایی انتخاب‌شده برای task قرار بده و مسیرهای شکل را مطابق همان نام اصلاح کن.

\subsubsection{ساختار داده و سیگنال‌ها}

برای این فعالیت، انتظار می‌رود حسگرهای پا نقش مهمی در ثبت تغییرات اصلی حرکت داشته باشند، زیرا بالا رفتن از جعبه نیازمند اعمال نیرو، جابه‌جایی وزن و کنترل وضعیت پا است. با این حال، حسگرهای دست و سر نیز می‌توانند الگوهای جبرانی، تغییرات تعادل و حرکت تنه را نشان دهند. بنابراین، تحلیل نهایی باید مشخص کند که کدام حسگرها و کدام محورهای حرکتی بیشترین نقش را در تفکیک کلاس‌ها داشته‌اند.

% TODO: شکل نمونه سیگنال فعالیت Box_Step را اینجا اضافه کن.

\subsubsection{تحلیل ویژگی‌ها و کاهش بعد}

پس از استخراج ویژگی‌ها، نقشه همبستگی و نمودارهای \lr{PCA} باید بررسی شوند. اگر فضای ویژگی دارای همبستگی زیاد باشد، کاهش بعد می‌تواند به ساده‌سازی مدل کمک کند. همچنین، پراکندگی داده‌ها در دو مؤلفه اول می‌تواند نشان دهد که آیا کلاس‌ها در این فعالیت از هم فاصله قابل مشاهده دارند یا خیر.

% TODO: شکل‌های feature_correlation، pca_variance و pca_scatter مربوط به Box_Step را اینجا اضافه کن.

\subsubsection{ارزیابی مدل‌ها}

برای این فعالیت نیز باید عملکرد سه مدل \lr{KNN}، \lr{SVM} و \lr{Random Forest} در حالت بدون \lr{PCA} و با \lr{PCA95} بررسی شود. اگر داده‌های جداسازی‌شده با تشخیص فعالیت برای این فعالیت تولید شوند، مقایسه آن‌ها با داده‌های اصلی نیز اهمیت خواهد داشت، زیرا بازه شروع و پایان بالا رفتن از جعبه می‌تواند بر ویژگی‌های آماری اثرگذار باشد.

% TODO: عددهای جدول زیر را پس از تولید Excel فعالیت Box_Step وارد کن.
\begin{table}[H]
	\centering
	\caption{خلاصه عملکرد مدل‌ها برای فعالیت بالا رفتن از جعبه}
	\label{tab:box_step_summary}
	\begin{tabular}{lccccc}
		\toprule
		مدل & \lr{Accuracy} & \lr{Precision} & \lr{Recall} & \lr{Specificity} & \lr{F1} \\
		\midrule
		\lr{KNN} &  &  &  &  &  \\
		\lr{SVM} &  &  &  &  &  \\
		\lr{Random Forest} &  &  &  &  &  \\
		\bottomrule
	\end{tabular}
\end{table}

\subsubsection{جمع‌بندی مورد انتظار}

پس از تکمیل داده‌ها، این فعالیت می‌تواند نقش مهمی در ارزیابی توان سامانه برای تحلیل حرکات عملکردی پیچیده‌تر داشته باشد. انتظار می‌رود ویژگی‌های مرتبط با اندام‌های تحتانی، تعادل و تغییرات تنه در این فعالیت مهم باشند. در تحلیل نهایی باید مشخص شود که آیا مدل منتخب برای این فعالیت با مدل‌های منتخب فعالیت‌های دیگر مشابه است یا به تنظیمات متفاوتی نیاز دارد.
